\section{Контрольная по занятиям №1--3}\label{answers:mock13}
\paragraph*{В1.}\label{ans:v1}
$205_{10}=11001101_2=CD_{16}$; $57_8+26_8=105_8$.
\hyperref[sol:v1]{Решение и баллы.}

\paragraph*{В2.}\label{ans:v2}
7-битный дополнительный код: $[-64,63]$, 128 значений. Для $-27$ в 8 битах: прямой \texttt{9B}, обратный \texttt{E4}, дополнительный \texttt{E5}, со смещением 128 --- \texttt{65}.
\hyperref[sol:v2]{Решение и баллы.}

\paragraph*{В3.}\label{ans:v3}
\begin{center}
\begin{tabular}{l|c|r|cccc}
Операция & Код & Signed & CF & OF & ZF & SF\\\hline
$85+60$ & \texttt{91} & $-111$ & 0 & 1 & 0 & 1\\
$40-75$ & \texttt{DD} & $-35$ & 1 & 0 & 0 & 1
\end{tabular}
\end{center}
\hyperref[sol:v3]{Решение и баллы.}

\paragraph*{В4.}\label{ans:v4}
AND даёт \texttt{10}; сброс бита 5 --- \texttt{93}; установка бита 2 --- \texttt{B7}; извлечённое поле --- \texttt{04}, число 4.
\hyperref[sol:v4]{Решение и баллы.}

\paragraph*{В5.}\label{ans:v5}
Логический вправо: \texttt{35} (53); арифметический вправо: \texttt{F5} ($-11$); циклический влево: \texttt{AB}; расширение: \texttt{FFD5}. Исходное знаковое число $-43$; деление на 4 с округлением к нулю даёт $-10$.
\hyperref[sol:v5]{Решение и баллы.}

\paragraph*{В6.}\label{ans:v6}
$-13.375=-1.101011_2\cdot2^3$. Поля \texttt{float32}:
\begin{center}
\texttt{1 10000010 10101100000000000000000}
\end{center}
Hex: \texttt{C1560000}. Для \texttt{40B80000}: знак 0, хранимый порядок 129, истинный порядок 2, значащая часть $1.0111_2$; значение $5.75$.
\hyperref[sol:v6]{Решение и баллы.}

\paragraph*{В7.}\label{ans:v7}
\texttt{80000000}: $-0$; \texttt{7F800000}: $+\infty$; \texttt{7FC00001}: NaN; \texttt{00400000}: положительное денормализованное число $2^{-127}$.
\hyperref[sol:v7]{Решение и баллы.}

\paragraph*{В8.}\label{ans:v8}
$\varepsilon=2^{-23}$. С округлением каждой операции: $(1+h)+h=1$, а $1+(h+h)=1+2^{-23}$. Коды результатов: \texttt{3F800000} и \texttt{3F800001}. Проявляются поглощение и потеря ассоциативности.
\hyperref[sol:v8]{Решение и баллы.}
